%% ch04_training_data.tex -- Intelligence Engineering, chapter 04 "Training
%% Data".
%% Spine transferred from Huyen, Designing Machine Learning Systems, chapter
%% 4. The subchapters and the frame titles under them are the book's own
%% section and subsection headings, in the book's order.
%%
%% STATE: titles and TEMPLATE layout only. No body text. Every frame carries
%% the layout it is meant to be filled into, and a % TEMPLATE comment saying
%% why that layout and what belongs in each slot. Filling them is the next pass.
%%
%% Fragment of intelligence_engineering.tex. One chapter is one lecture and
%% gets exactly ONE \lecturepage, at the top. Inside it every \subsection
%% opens on the subchapter divider, which the master emits automatically via
%% \AtBeginSubsection; do not write those divider frames by hand.

\begin{refsection}

\lecturepage{IV}{Training Data}{}
\section[Training Data]{Training Data}

% ----------------------------------------------------------------------
\subsection{Sampling}

% TEMPLATE: withmargin carrying a definition box. The heading is the term the
% first half of this subchapter turns on, so it takes the slide's one titled
% box and the book's kinds of nonprobability sampling are listed inside it.
% Margin takes the aside the book uses to justify opening here, the convenience
% of these samples against how unrepresentative they are.
\begin{frame}{Nonprobability Sampling}
  \begin{withmargin}
    \begin{tdmain}
      \begin{tddef}{Nonprobability sampling}
      \end{tddef}
    \end{tdmain}
    \begin{tdmargin}
    \end{tdmargin}
  \end{withmargin}
  \slidesources{Huyen2022}
\end{frame}

% TEMPLATE: two levelled boxes, untitled on purpose. The book gives this method
% one selection rule and one failure mode, which read as two equal halves
% rather than as a claim with an aside. The left box takes the rule and the
% probability it assigns each sample, the right box what it does to a rare
% class; the later pass opens each box with its own bold label.
\begin{frame}{Simple Random Sampling}
  \begin{columns}[T,onlytextwidth]
    \begin{column}{0.47\textwidth}
      \begin{tdblock}[equal height group=srs]
      \end{tdblock}
    \end{column}
    \begin{column}{0.47\textwidth}
      \begin{tdblock}[equal height group=srs]
      \end{tdblock}
    \end{column}
  \end{columns}
  \slidesources{Huyen2022}
\end{frame}

% TEMPLATE: withmargin. One mechanism plus one caveat, which is exactly the
% shape this layout is for. Main column takes how the population is divided
% into groups and sampled inside each of them, with the book's own class
% proportions. Margin takes the limit the book immediately raises, that the
% groups cannot always be drawn.
\begin{frame}{Stratified Sampling}
  \begin{withmargin}
    \begin{tdmain}
    \end{tdmain}
    \begin{tdmargin}
    \end{tdmargin}
  \end{withmargin}
  \slidesources{Huyen2022}
\end{frame}

% TEMPLATE: code panel captioned with the heading. The book makes this method
% concrete with a weights vector and a single call rather than with prose, and
% a weights vector only reads as itself in mono, so the panel is the whole
% slide. It takes that snippet, the weights it passes and the distribution they
% encode.
\begin{frame}{Weighted Sampling}
  \begin{tdcodet}{Weighted sampling}
  \end{tdcodet}
  \slidesources{Huyen2022}
\end{frame}

% TEMPLATE: three levelled boxes. The book states this one as an algorithm of
% three moves over a stream, and the moves only make sense in sequence, so one
% box each rather than a paragraph. The equal height group keeps the row level;
% the boxes stay untitled until the later pass numbers them and adds the
% probability that makes the reservoir uniform.
\begin{frame}{Reservoir Sampling}
  \begin{columns}[T,onlytextwidth]
    \begin{column}{0.31\textwidth}
      \begin{tdblock}[equal height group=resv]
      \end{tdblock}
    \end{column}
    \begin{column}{0.31\textwidth}
      \begin{tdblock}[equal height group=resv]
      \end{tdblock}
    \end{column}
    \begin{column}{0.31\textwidth}
      \begin{tdblock}[equal height group=resv]
      \end{tdblock}
    \end{column}
  \end{columns}
  \slidesources{Huyen2022}
\end{frame}

% TEMPLATE: withmargin. The book's case rests on one ratio between two
% distributions, so the main column carries that expression and the condition
% on the proposal distribution. Margin takes the use the book reaches for,
% where the samples are expensive and the distribution you want is not the one
% you can draw from.
\begin{frame}{Importance Sampling}
  \begin{withmargin}
    \begin{tdmain}
    \end{tdmain}
    \begin{tdmargin}
    \end{tdmargin}
  \end{withmargin}
  \slidesources{Huyen2022}
\end{frame}

% ----------------------------------------------------------------------
\subsection{Labeling}

% TEMPLATE: three levelled boxes, and the opener for the two frames under this
% heading. The book's case against hand labels is three costs of the same kind,
% peers of one another, so one box each rather than a bullet list. The two
% named problems below get their own frames and are deliberately not repeated
% here; the boxes stay untitled until the later pass names each cost.
\begin{frame}{Hand Labels}
  \begin{columns}[T,onlytextwidth]
    \begin{column}{0.31\textwidth}
      \begin{tdblock}[equal height group=hand]
      \end{tdblock}
    \end{column}
    \begin{column}{0.31\textwidth}
      \begin{tdblock}[equal height group=hand]
      \end{tdblock}
    \end{column}
    \begin{column}{0.31\textwidth}
      \begin{tdblock}[equal height group=hand]
      \end{tdblock}
    \end{column}
  \end{columns}
  \slidesources{Huyen2022}
\end{frame}

% TEMPLATE: withmargin. The book argues this one from a single annotated
% sentence and the disagreement it produces, so the main column takes that
% example and the differing counts it yields. Margin takes the book's remedy,
% the problem definition and the instructions that have to exist before the
% annotators start.
\begin{frame}{Label multiplicity}
  \begin{withmargin}
    \begin{tdmain}
    \end{tdmain}
    \begin{tdmargin}
    \end{tdmargin}
  \end{withmargin}
  \slidesources{Huyen2022}
\end{frame}

% TEMPLATE: definition box alone, full width. The heading is the book's own
% term and there is nothing here to contrast it against, so the definition
% takes the whole frame instead of sharing it. The box takes what has to be
% tracked per sample and the failure the book reports when it is not, a model
% that got worse after new data arrived.
\begin{frame}{Data lineage}
  \begin{tddef}{Data lineage}
  \end{tddef}
  \slidesources{Huyen2022}
\end{frame}

% TEMPLATE: booktabs table, pushed there on purpose rather than left as a
% withmargin. The book names its canonical cases in prose, but every one of
% them is the same two fields, the task and the label the system gets for free,
% so a table is the tighter transfer and makes the pattern visible. One row per
% case, column headers taken from the heading.
\begin{frame}{Natural Labels}
  \footnotesize
  \renewcommand{\arraystretch}{1.25}
  \begin{tabular}{@{}>{\raggedright\arraybackslash}p{0.38\textwidth}>{\raggedright\arraybackslash}p{0.50\textwidth}@{}}
    \toprule
    \textbf{Task} & \textbf{Natural label}\\
    \midrule
     & \\
     & \\
     & \\
     & \\
    \bottomrule
  \end{tabular}
  \slidesources{Huyen2022}
\end{frame}

% TEMPLATE: two levelled boxes. The heading is about a length, and a length
% only means something as a short end read against a long end, so each end
% takes one box headed by it. Each box takes the book's own examples and the
% window the label closes in, which is the number the deployment decision
% later hangs on.
\begin{frame}{Feedback loop length}
  \begin{columns}[T,onlytextwidth]
    \begin{column}{0.47\textwidth}
      \begin{tdblockt}[equal height group=fbl]{Short feedback loops}
      \end{tdblockt}
    \end{column}
    \begin{column}{0.47\textwidth}
      \begin{tdblockt}[equal height group=fbl]{Long feedback loops}
      \end{tdblockt}
    \end{column}
  \end{columns}
  \slidesources{Huyen2022}
\end{frame}

% TEMPLATE: two by two block grid, and the opener for the four frames under
% this heading. The heading introduces exactly four peer methods, none
% subordinate to another, so each takes one box headed with its own name rather
% than one long list. Each box carries only the one line the book's comparison
% uses to tell that method apart, how it gets its labels and whether it needs
% ground truths; the detail waits for the frames below.
\begin{frame}{Handling the Lack of Labels}
  \begin{columns}[T,onlytextwidth]
    \begin{column}{0.47\textwidth}
      \begin{tdblockt}[equal height group=lack]{Weak supervision}
      \end{tdblockt}
      \vskip0.5em
      \begin{tdblockt}[equal height group=lack2]{Transfer learning}
      \end{tdblockt}
    \end{column}
    \begin{column}{0.47\textwidth}
      \begin{tdblockt}[equal height group=lack]{Semi-supervision}
      \end{tdblockt}
      \vskip0.5em
      \begin{tdblockt}[equal height group=lack2]{Active learning}
      \end{tdblockt}
    \end{column}
  \end{columns}
  \slidesources{Huyen2022}
\end{frame}

% TEMPLATE: side image cover, the first of this chapter's two image slots. The
% book's own content here is a pipeline, heuristics written as functions that
% vote and are combined into labels, and a pipeline is a figure before it is a
% sentence, so the strip holds it. The body holds the definition of a labeling
% function and the kinds of heuristic the book lists.
\begin{frame}{Weak supervision}
  \phimgside[0.33]{weak_supervision.png}{Weak supervision}
  \begin{sidebody}
  \end{sidebody}
  \slidesources{Huyen2022}
\end{frame}

% TEMPLATE: three levelled boxes. The book names three ways to grow labels from
% a small seed set, peers of one another, so one box each rather than a bullet
% list, levelled by the equal height group. The boxes stay untitled until the
% later pass opens each with its method name and the assumption it rests on.
\begin{frame}{Semi-supervision}
  \begin{columns}[T,onlytextwidth]
    \begin{column}{0.31\textwidth}
      \begin{tdblock}[equal height group=semi]
      \end{tdblock}
    \end{column}
    \begin{column}{0.31\textwidth}
      \begin{tdblock}[equal height group=semi]
      \end{tdblock}
    \end{column}
    \begin{column}{0.31\textwidth}
      \begin{tdblock}[equal height group=semi]
      \end{tdblock}
    \end{column}
  \end{columns}
  \slidesources{Huyen2022}
\end{frame}

% TEMPLATE: withmargin. One mechanism plus one aside. Main column takes the
% book's order, the base task and its data, then the fine tuning step, then the
% zero shot case where no fine tuning happens at all. Margin takes what the
% book says this did to the cost of entry for tasks with little data.
\begin{frame}{Transfer learning}
  \begin{withmargin}
    \begin{tdmain}
    \end{tdmain}
    \begin{tdmargin}
    \end{tdmargin}
  \end{withmargin}
  \slidesources{Huyen2022}
\end{frame}

% TEMPLATE: two levelled boxes. The book's heuristics for choosing which sample
% to label next fall into two families, so they read side by side rather than
% as one list. Each box takes one family, the quantity it ranks candidates by
% and the book's query strategy for it; the later pass heads each box with the
% family name.
\begin{frame}{Active learning}
  \begin{columns}[T,onlytextwidth]
    \begin{column}{0.47\textwidth}
      \begin{tdblock}[equal height group=actv]
      \end{tdblock}
    \end{column}
    \begin{column}{0.47\textwidth}
      \begin{tdblock}[equal height group=actv]
      \end{tdblock}
    \end{column}
  \end{columns}
  \slidesources{Huyen2022}
\end{frame}

% ----------------------------------------------------------------------
\subsection{Class Imbalance}

% TEMPLATE: withmargin plus a takeaway. The book argues from one ratio and then
% from three consequences of it, so the main column runs them in that order and
% the margin carries the worked case it uses, the scan set and its counts. The
% takeaway keeps the book's punch in front of the claim, the model that scores
% well and never predicts the rare class.
\begin{frame}{Challenges of Class Imbalance}
  \begin{withmargin}
    \begin{tdmain}
    \end{tdmain}
    \begin{tdmargin}
    \end{tdmargin}
  \end{withmargin}
  \begin{takeaway}
  \end{takeaway}
  \slidesources{Huyen2022}
\end{frame}

% TEMPLATE: three levelled boxes, and the opener for the three frames under
% this heading. The heading introduces three peer families of fix and the book
% keeps them strictly apart, so one box each headed with its own name. Each box
% holds only the one line that tells that family from the other two, the level
% it acts on; the detail waits for the frames below.
\begin{frame}{Handling Class Imbalance}
  \begin{columns}[T,onlytextwidth]
    \begin{column}{0.31\textwidth}
      \begin{tdblockt}[equal height group=hci]{Evaluation metrics}
      \end{tdblockt}
    \end{column}
    \begin{column}{0.31\textwidth}
      \begin{tdblockt}[equal height group=hci]{Data-level methods}
      \end{tdblockt}
    \end{column}
    \begin{column}{0.31\textwidth}
      \begin{tdblockt}[equal height group=hci]{Algorithm-level methods}
      \end{tdblockt}
    \end{column}
  \end{columns}
  \slidesources{Huyen2022}
\end{frame}

% TEMPLATE: booktabs table with the header row left open. The book makes this
% point with a model by metric grid in which the same two models swap ranks the
% moment the metric changes, so the grid is the argument and a paragraph would
% lose it. The header row takes the book's metric names, the rows its two
% models, and the cells the numbers that reverse.
\begin{frame}{Using the right evaluation metrics}
  \footnotesize
  \renewcommand{\arraystretch}{1.25}
  \begin{tabular}{@{}llll@{}}
    \toprule
     & & & \\
    \midrule
     & & & \\
     & & & \\
    \bottomrule
  \end{tabular}
  \slidesources{Huyen2022}
\end{frame}

% TEMPLATE: two levelled boxes. Resampling runs in two directions and the
% book's named methods split cleanly along them, so one box each rather than
% one list that hides the split. Each box takes its direction, the methods the
% book names for it and what that direction risks, losing data on one side and
% overfitting on the other.
\begin{frame}{Data-level methods: Resampling}
  \begin{columns}[T,onlytextwidth]
    \begin{column}{0.47\textwidth}
      \begin{tdblock}[equal height group=rsmp]
      \end{tdblock}
    \end{column}
    \begin{column}{0.47\textwidth}
      \begin{tdblock}[equal height group=rsmp]
      \end{tdblock}
    \end{column}
  \end{columns}
  \slidesources{Huyen2022}
\end{frame}

% TEMPLATE: withmargin. The book's three modifications are all the same move on
% the loss function, so the main column carries them together with their
% expressions rather than splitting them into boxes that would hide how alike
% they are. Margin carries the shared idea, that a mistake on the rare class
% has to cost more than a mistake on the common one.
\begin{frame}{Algorithm-level methods}
  \begin{withmargin}
    \begin{tdmain}
    \end{tdmain}
    \begin{tdmargin}
    \end{tdmargin}
  \end{withmargin}
  \slidesources{Huyen2022}
\end{frame}

% ----------------------------------------------------------------------
\subsection{Data Augmentation}

% TEMPLATE: two levelled boxes. The book gives this one twice, once for images
% and once for text, the same move in two modalities, so the two modalities sit
% side by side rather than in one list. Each box takes that modality's
% transformations in the book's words and the reason the label survives them.
\begin{frame}{Simple Label-Preserving Transformations}
  \begin{columns}[T,onlytextwidth]
    \begin{column}{0.47\textwidth}
      \begin{tdblock}[equal height group=slpt]
      \end{tdblock}
    \end{column}
    \begin{column}{0.47\textwidth}
      \begin{tdblock}[equal height group=slpt]
      \end{tdblock}
    \end{column}
  \end{columns}
  \slidesources{Huyen2022}
\end{frame}

% TEMPLATE: side image cover, this chapter's second and last image slot. The
% book makes this point with a picture, a sample a human still reads one way
% while the model reads it another after a change too small to see, and that
% only lands as an image. The strip holds it; the body holds the definition and
% the defensive use the book draws from it.
\begin{frame}{Perturbation}
  \phimgside[0.33]{perturbation.png}{Perturbation}
  \begin{sidebody}
  \end{sidebody}
  \slidesources{Huyen2022}
\end{frame}

% TEMPLATE: withmargin carrying a code panel. The book's example here is a
% template with slots that get filled to generate samples, and a template only
% reads as itself in mono, so the panel takes it in the main column. Margin
% takes the boundary the book draws, what synthesis can stand in for and what
% it cannot.
\begin{frame}{Data Synthesis}
  \begin{withmargin}
    \begin{tdmain}
      \begin{tdcode}
      \end{tdcode}
    \end{tdmain}
    \begin{tdmargin}
    \end{tdmargin}
  \end{withmargin}
  \slidesources{Huyen2022}
\end{frame}

% ----------------------------------------------------------------------
\subsection{Summary}

% TEMPLATE: bare takeaway. A summary slide that repeats the chapter as bullets
% is dead weight, and the book's own summary is one claim rather than a recap,
% so this frame carries a single centred line and nothing else, the book's
% closing statement on training data and the work it still demands.
\begin{frame}{Summary}
  \begin{takeaway}
  \end{takeaway}
  \slidesources{Huyen2022}
\end{frame}

% ----------------------------------------------------------------------
% This chapter's own references. The refsection opened above scopes the
% citation numbering to this chapter, so chapter_pdfs/<this>.pdf resolves
% its own [n] markers instead of pointing at a list it does not contain.
\begin{frame}{References}
  \printbibliography[heading=none]
\end{frame}
\end{refsection}
